% Part: computability
% Chapter: computability-theory
% Section: complement-ce

\documentclass[../../../include/open-logic-section]{subfiles}

\begin{document}

\olfileid{cmp}{thy}{cmp}
\olsection{المجموعات القابلة للتعداد حسابيًا ليست منغلقة تحت المتممة}

إذا كانت ${\OLId{latin-looped}{A}}$ قابلة للتعداد حسابيًا، أيلزم أن تكون متممة~${\OLId{latin-looped}{A}}$، أعني
$\Complement{{\OLId{latin-looped}{A}}}=\Nat\setminus {\OLId{latin-looped}{A}}$، قابلة للتعداد حسابيًا؟
ليس ذلك بلازم، كما تبيّن المبرهنة التالية ونتيجتها.

\begin{thm}
\ollabel{thm:ce-comp}
لكل مجموعة أعداد طبيعية ${\OLId{latin-looped}{A}}$، تتكافأ قابلية ${\OLId{latin-looped}{A}}$ للحساب
مع كون ${\OLId{latin-looped}{A}}$ و$\Complement{{\OLId{latin-looped}{A}}}$ جميعًا قابلتين للتعداد
حسابيًا.
\end{thm}

\begin{proof}
أما الوجه الأول فظاهر؛ فإن قابلية ${\OLId{latin-looped}{A}}$ للحساب تستلزم قابلية $\Complement{{\OLId{latin-looped}{A}}}$
له ($\Char{{\OLId{latin-looped}{A}}}={\OLMathNumeral{1}}\tsub\Char{\Complement{{\OLId{latin-looped}{A}}}}$)، ومن ثم
تكونان كلتاهما قابلتين للتعداد حسابيًا.

وأما العكس، فلنفرض أن ${\OLId{latin-looped}{A}}$ و~$\Complement{{\OLId{latin-looped}{A}}}$ قابلتان
للتعداد حسابيًا. نجعل ${\OLId{latin-looped}{A}}$ مجال~$\cfind{{\OLId{latin-ordinary}{d}}}$،
و$\Complement{{\OLId{latin-looped}{A}}}$ مجال~$\cfind{{\OLId{latin-ordinary}{e}}}$، ونعرّف ${\OLId{latin-ordinary}{h}}$ كما يأتي:
\[
{\OLId{latin-ordinary}{h}}({\OLId{latin-ordinary}{x}}) = \umin{{\OLId{latin-ordinary}{s}}}{({\OLId{latin-looped}{T}}({\OLId{latin-ordinary}{d}},{\OLId{latin-ordinary}{x}},{\OLId{latin-ordinary}{s}}) \lor {\OLId{latin-looped}{T}}({\OLId{latin-ordinary}{e}},{\OLId{latin-ordinary}{x}},{\OLId{latin-ordinary}{s}}))}.
\]
فالدالة~${\OLId{latin-ordinary}{h}}$ تبحث عند الدخل~${\OLId{latin-ordinary}{x}}$ عن سجل حساب متوقف لـ
~$\cfind{{\OLId{latin-ordinary}{d}}}$ أو عن سجل حساب متوقف لـ~$\cfind{{\OLId{latin-ordinary}{e}}}$. فإذا كان ${\OLId{latin-ordinary}{x}}\in {\OLId{latin-looped}{A}}$،
% تصحيح صياغة كلاسيكية (C242-01): وُحّد حرف الجر في معمولي البحث المتوازيين.
نجحت في الحالة الأولى، وإذا كان ${\OLId{latin-ordinary}{x}}\in\Complement{{\OLId{latin-looped}{A}}}$، نجحت في الحالة
الثانية. فلا بد من نجاح البحث؛ وتكون~${\OLId{latin-ordinary}{h}}$ كلية قابلة للحساب. ثم إن كل
~${\OLId{latin-ordinary}{x}}$ يحقق ${\OLId{latin-ordinary}{x}}\in {\OLId{latin-looped}{A}}$ إذا وفقط إذا صدقت ${\OLId{latin-looped}{T}}({\OLId{latin-ordinary}{d}},{\OLId{latin-ordinary}{x}},{\OLId{latin-ordinary}{h}}({\OLId{latin-ordinary}{x}}))$، أي إذا كانت
$\cfind{{\OLId{latin-ordinary}{d}}}$ هي المعرفة. وحيث تقبل العلاقة ${\OLId{latin-looped}{T}}({\OLId{latin-ordinary}{d}},{\OLId{latin-ordinary}{x}},{\OLId{latin-ordinary}{h}}({\OLId{latin-ordinary}{x}}))$ الحساب،
تقبل~${\OLId{latin-looped}{A}}$ الحساب أيضًا.
\end{proof}

\begin{explain}
ويقرب ذلك بوصف إجرائي: لقرار ${\OLId{latin-looped}{A}}$ نبحث عند الدخل
${\OLId{latin-ordinary}{x}}$ عن سجل توقف لـ$\cfind{{\OLId{latin-ordinary}{d}}}$ أو $\cfind{{\OLId{latin-ordinary}{e}}}$. ولا يفوتنا
أحدهما؛ فإن توقف $\cfind{{\OLId{latin-ordinary}{d}}}$، كان ${\OLId{latin-ordinary}{x}}$ من~${\OLId{latin-looped}{A}}$، وإلا كان ${\OLId{latin-ordinary}{x}}$ من
~$\Complement{{\OLId{latin-looped}{A}}}$.
\end{explain}

\begin{cor}
\ollabel{cor:comp-k}
$\Complement{{\OLId{latin-looped}{K}}_{\OLMathNumeral{0}}}$ غير قابلة للتعداد حسابيًا.
\end{cor}

\begin{proof}
قد علمنا قابلية ${\OLId{latin-looped}{K}}_{\OLMathNumeral{0}}$ للتعداد حسابيًا مع عدم قابليتها للحساب. ولو كانت
$\Complement{{\OLId{latin-looped}{K}}_{\OLMathNumeral{0}}}$ قابلة للتعداد حسابيًا، لكانت ${\OLId{latin-looped}{K}}_{\OLMathNumeral{0}}$ قابلة للحساب
بحسب \olref{thm:ce-comp}، خلافًا لـ\olref[ncp]{thm:K-0}.
\end{proof}

\end{document}
